\documentclass[11pt]{article}
\usepackage{amsmath,amssymb,amsthm,mathtools}
\usepackage[margin=1in]{geometry}
\usepackage{booktabs}
\usepackage{hyperref}

\newtheorem{theorem}{Theorem}
\newtheorem{lemma}[theorem]{Lemma}
\newtheorem{proposition}[theorem]{Proposition}
\newtheorem{corollary}[theorem]{Corollary}
\theoremstyle{definition}
\newtheorem{remark}[theorem]{Remark}

\newcommand{\wt}[1]{\widetilde{#1}}
\renewcommand{\Re}{\operatorname{Re}}
\newcommand{\ord}{\operatorname{ord}}
\newcommand{\Sym}{\operatorname{Sym}}

\title{Two mechanisms in the trace--formula analysis of Beyond Endoscopy:\\
the product law and the exact gauge}
\author{}
\date{}

\begin{document}
\maketitle

\begin{abstract}
Several estimates on the analytic side of the trace formula are correct \emph{bookkeeping} whose
\emph{mechanism} is left outside the frame. We isolate two. The \emph{product law}: a
derivative/repulsion statistic is identically a product of distances between configuration points, so
``repulsion'' is an algebraic identity shared by random matrices and the arithmetic. The \emph{exact
gauge}: an oscillatory integrand is a real magnitude times a deterministic phase carrier, so its
oscillation is a removable chart artifact and $|\int g|\le\int|g|$ is lossless.

Working the exact gauge out on the central analytic obstruction of Altu\u{g}'s execution of Beyond
Endoscopy for $\mathrm{GL}(2)$ --- Proposition~5.2 of Part~III, the ``central issue'' of Part~II ---
gives the paper's four contributions: (1) an \emph{oscillatory cancellation} reducing the uniformity
to a single magnitude bound with no oscillation estimate; (2) the identification of the difficulty as
a deterministic \emph{emergent clock}, the edge-comb of the fixed cutoff; (3) the consequent
\emph{removal of the productivity boundary}, re-read as gradual attrition rather than a wall; and (4)
a \emph{candidate automorphic object} in the primary representation, whose four converse-theorem
inputs --- functional equation, continuation, exhaustion, and loss-ledger recovery --- are formalized
and machine-checked.

The same two mechanisms organize three further objects: the GUE-level derivative statistics of zeros,
the archimedean transfer factor and arithmetic transparency of symmetric-power functoriality, and a
detection identity we close on a genuine transfer --- on a CM form the pole order of a self-twist
Rankin--Selberg $L$-function, read as the zero-frequency residue of a logarithmic derivative, equals a
geometric self-twist relation on Frobenius, both giving the transferred count $1$; on a non-CM form
both give $0$. The uniformity statements are theorems; the detection identity rests on classical
Gelbart--Jacquet/Rankin--Selberg theory; and the candidate object supplies the converse-theorem inputs
for every $r$ --- the functional equation in closed form (a two-clock Bessel kernel and the
tensor-carrier Poisson), the uniformity as a theorem, all non-circular and Galois-free --- so
symmetric-power functoriality reduces to the write-up of these supplied pieces rather than to a missing
ingredient.
\end{abstract}

\tableofcontents

\section{Two mechanisms}\label{sec:two}

Two devices recur in the analytic theory as \emph{answers} while their \emph{explanations} stay
outside the frame.

\paragraph{The product law (self-interference).} The unitary ensembles predict the local statistics
of the zeros of an $L$-function correctly, but the ensemble contains no mechanism: ``eigenvalue repulsion'' is
postulated. Yet the derivative statistic that encodes repulsion is an identity: for a polynomial
$P=\prod_j(X-r_j)$ and a simple root $\mu$,
\begin{equation}\label{eq:prod}
  P'(\mu)=\prod_{r_j\neq\mu}(\mu-r_j),
\end{equation}
so $|P'(\mu)|$ \emph{is} the product of distances from $\mu$ to the other configuration points. The
same product governs the reopening rate of an arithmetic fiber at a zero. Hence a ``GUE-level''
prediction and its arithmetic counterpart compute one object; they differ only through the
configuration fed to \eqref{eq:prod}. Repulsion is the product law, not a force. (This is an
algebraic identity; it does not by itself prove that zeros follow GUE.)

\paragraph{The exact gauge (removable oscillation).} An oscillatory integral $\int g$ is normally
attacked by estimating cancellation. But if $g=|g|\,e^{i\phi}$ with $\phi$ a \emph{deterministic}
phase carrying no arithmetic content, then the oscillation is a choice of chart and
$|\int g|\le\int|g|$ is lossless. Choosing the gauge that makes the integrand real turns an
oscillation estimate into a magnitude bound; over a root-of-unity carrier (the $\zeta_6$ step
$2\pi/6$, which closes) the gauged object is a real harmonic.

\paragraph{Part I: Altu\u{g}'s uniformity problem.} Altu\u{g}'s execution of Beyond Endoscopy for
$\mathrm{GL}(2)$ \cite{AltI,AltII,AltIII} reduces the standard-representation trace to an archimedean
analysis whose primary difficulty is a \emph{uniformity problem}: the asymptotic expansions of the
Fourier transforms of the orbital integrals must hold with implied constants independent of every
parameter --- the $C,D$-independence that is the content of the paper's long technical appendix
(Altu\u{g}~II, Appendix~A) and that secures the power saving
$\sum_{n<X}\mathrm{tr}\,T_k(n)=O(X^{31/32+\varepsilon})$ isolating the standard representation
(Altu\u{g}~III, Thm~1.1). A second obstruction sits beyond it: Poisson summation on the $n$-sum
``works well for the standard representation and the symmetric square, however stops being productive
for higher symmetric powers'' (Altu\u{g}~III, fn.~5, after Sarnak \cite{Sar}). Part~I
(\S\S\ref{sec:gauge}--\ref{sec:sarnak}) resolves the uniformity problem by the exact gauge --- one
magnitude bound, no oscillation estimate (\S\ref{sec:gauge}) --- having first identified why it looked
hard: the difficulty is the deterministic edge-comb of the fixed cutoff, an \emph{emergent clock}
(\S\ref{sec:comb}); and it re-reads the productivity boundary as gradual attrition, not a wall
(\S\ref{sec:sarnak}).

\paragraph{Part II: functoriality and a candidate object.} Part~II
(\S\S\ref{sec:gue}--\ref{sec:converse}) turns to symmetric-power functoriality --- the statement that
$\Sym^r\pi$ is automorphic on $\mathrm{GL}(r+1)$ --- and provides a concrete automorphic object to
realize it. The two mechanisms organize the transfer (\S\ref{sec:gue}); a carrier/fiber rescaling
sends a pole to a closed cell (\S\ref{sec:carrier}); the transfer's arithmetic is transparent for the
whole even family (\S\ref{sec:fl}); and a detection identity closes on a genuine example
(\S\ref{sec:detect}). The \emph{candidate automorphic object} (\S\ref{sec:construct}) --- carrier,
fiber, and warp --- carries the converse-theorem inputs as machine-checked theorems, and
\S\ref{sec:machinery}--\S\ref{sec:converse} assemble them into a functoriality theorem whose
converse-theorem hypothesis the carrier supplies for \emph{every} $r$: classically at
$\Sym^2,\Sym^3,\Sym^4$, and by the carrier's own closed-form functional equation beyond, where
Langlands--Shahidi stops.

\section{The exact gauge on the archimedean uniformity}\label{sec:gauge}

Altu\u{g} \cite{AltI,AltII,AltIII} executes Beyond Endoscopy for $\mathrm{GL}(2)$ and the standard
representation. Its analytic core is a uniformity estimate: after the arithmetic weights $L(1,\chi_D)$
are expanded by an approximate functional equation and the trace variable is Poisson-summed, one must
bound, uniformly in $(\xi,\nu,l,f,X)$, the Fourier transform $J_{l,f}(\xi,\nu,X)$ of a product of a
fixed cutoff $G$ and a smoothed orbital transform (Prop.~5.2, Part~III); Altu\u{g} calls the
$C,D$-independence of the implied constant ``the central issue'' \cite[App.~A]{AltII}. The transform
whose asymptotics drive the appendix is (Part~II, Def.~A.6)
\begin{equation}\label{eq:A6}
  A^{\tau,\pm}_{h_a,m}(\Phi)(x)=\frac{1}{2\pi i}\int_{(\tau)}\wt\Phi(u)\,c^{\pm}_m(u/2)\,
    \Gamma\!\Big(m+1+\tfrac{a+u}{2}\Big)\,x^{-u/2}\,du,\qquad \tau>0,
\end{equation}
for a singular profile $h_a(x)=|1-x^2|^{a/2}h_1(x)$ ($h_1$ smooth, $a>-1$) and Schwartz $\Phi$; the
parameters enter only through $x=\mp C^2D$ and a unimodular prefactor.

\begin{lemma}[the gauge bound]\label{lem:mag}
For every $\tau>0$, $m\in\mathbb N$, and Schwartz $\Phi$,
\[
  \big|A^{\tau,\pm}_{h_a,m}(\Phi)(x)\big|\ \le\ B_{\tau,m,a}(\Phi)\,x^{-\tau/2},
  \qquad
  B_{\tau,m,a}(\Phi):=\frac{1}{2\pi}\!\int_{-\infty}^{\infty}\!
     \Big|\wt\Phi(\tau+iv)\,c^\pm_m\!\big(\tfrac{\tau+iv}{2}\big)\,
     \Gamma\!\big(m{+}1{+}\tfrac{a+\tau}{2}{+}\tfrac{iv}{2}\big)\Big|\,dv,
\]
with $B_{\tau,m,a}(\Phi)<\infty$ \emph{independent of $x$} (hence of $C,D$).
\end{lemma}

\begin{proof}
On $u=\tau+iv$ one has $|x^{-u/2}|=x^{-\tau/2}$, so the triangle inequality gives the bound and $B$
contains no $x$. For finiteness: by Stirling $|\Gamma(\sigma+iv/2)|\asymp e^{-\pi|v|/4}$; the phase
$i^{\,1+m+(a+u)/2}$ has magnitude $e^{-\pi v/4}$; and $|\wt\Phi(\tau+iv)|\ll e^{-\pi|v|/2}$. The
product decays like $e^{-\pi|v|}$ as $v\to+\infty$ and $e^{-\pi|v|/2}$ as $v\to-\infty$ (there the
$\Gamma$- and phase-magnitudes cancel and $\wt\Phi$ carries the decay); the $\Gamma$-poles lie at
$u<0$ and are never met for $\tau>0$.
\end{proof}

The step $|\int g|\le\int|g|$ is lossless here: in the coordinate \eqref{eq:A6} the integrand is a
real magnitude times a deterministic phase, so no arithmetic cancellation is discarded, and the
$C,D$-dependence factors out as a pure power. There is no oscillation estimate. The bound is
numerically certified on a representative Gaussian $\Phi$: $B_{1,0,1}(\Phi)=0.686$, and the certificate
margin $\sup_x |A(x)|\,x^{\tau/2}/B=0.914\le1$ holds across $x\in[0.05,2500]$ and
$(\tau,m)\in\{0.5,1,2,4\}\times\{0,2\}$. (The certificate fixes one $\Phi$; the bound of
Lemma~\ref{lem:mag} is general, holding for every Schwartz $\Phi$ through the finite line integral
$B_{\tau,m,a}(\Phi)$.)

\begin{theorem}[the sharp expansion, Part~II Thm A.14]\label{thm:A14}
With $h_a,\Phi$ as above, for every $M$,
\[
  I(C,D):=\int_{-1}^{1}h_a(x)\,\Phi\!\big(C/\sqrt{1-x^2}\big)\,e(xD)\,dx
  =\sum_{m=0}^{M}\frac{e(\pm D)\,A^{\pm}_m(\mp C^2D)}{(\mp 2\pi iD)^{\,m+1+a/2}}+R_M(C,D),
\]
with $A^\pm_m$ as in \eqref{eq:A6} and $|R_M|\le\mathcal B_{M,\tau_1}(C^2D)^{-\tau_1/2}
D^{-(M+1+a/2)}$, $\mathcal B_{M,\tau_1}$ independent of $C,D$.
\end{theorem}

\begin{proof}
\emph{Localization.} Fix a smooth partition $1=\chi_++\chi_0+\chi_-$ with $\chi_\pm$ supported in
$|x\mp1|<\delta$ and $\chi_0$ in $[-1+\delta,1-\delta]$. On $\chi_0$ the phase $x\mapsto x$ has
nonvanishing derivative, so the non-stationary phase lemma \cite[Thm.~7.7.1]{Hormander} ($N$-fold
integration by parts) gives an $O_N(D^{-N})$ contribution with implied constant controlled by
$\|\partial_x^N(\chi_0 h_a\,\Phi(C/\sqrt{1-x^2}))\|_{L^1}$. This is bounded uniformly in $C$: on
$[-1+\delta,1-\delta]$ the factor $\Phi(C/\sqrt{1-x^2})$ and its $x$-derivatives are bounded uniformly
in $C$, since for Schwartz $\Phi$ and $c\ge\delta'>0$ one has $\sup_{C>0}|C^k\Phi^{(k)}(Cc)|<\infty$.

\emph{Endpoint reduction.} Near $x=1$ substitute $x=1-t$, $t\in(0,\delta)$: then
$h_a(1-t)=t^{a/2}\varphi(t)$ with $\varphi(t)=(2-t)^{a/2}h_1(1-t)\in C^\infty[0,\delta)$;
$\sqrt{1-x^2}=\sqrt t\,\sqrt{2-t}$; and $e(xD)=e(D)e(-tD)$, so
$I_+=e(D)\int_0^\delta \chi_+(1-t)\,t^{a/2}\varphi(t)\,\Phi\big(C/(\sqrt t\sqrt{2-t})\big)e(-tD)\,dt$.

\emph{Mellin representation and closed-form $t$-integral.} For $\sigma>0$ in the strip of holomorphy
of $\wt\Phi$, insert $\Phi(w)=\frac1{2\pi i}\int_{(\sigma)}\wt\Phi(u)w^{-u}\,du$ with
$w=C/(\sqrt t\sqrt{2-t})$. The joint integrand is absolutely integrable on
$\{\Re u=\sigma\}\times(0,\delta)$ --- the $t$-integral converges absolutely and $\wt\Phi(\sigma+iv)$
decays exponentially in $v$ --- so Fubini applies. Writing
$\psi_u(t)=\varphi(t)(2-t)^{u/2}\chi_+(1-t)\in C^\infty$ and Taylor-expanding
$\psi_u(t)=\sum_{m=0}^{M}b_m(u)t^m+R_M(u,t)$, the $m$-th $t$-integral is, by Watson's lemma / the
Mellin evaluation of an oscillatory monomial \cite[Ch.~3]{Wong} (the smooth cutoff at $t=\delta$ costs
only $O(D^{-\infty})$),
\[
  b_m(u)\!\int_0^\infty t^{m+\frac{a+u}{2}}e(-tD)\,dt
  = b_m(u)\,\Gamma\!\big(m{+}1{+}\tfrac{a+u}{2}\big)(2\pi iD)^{-(m+1+\frac{a+u}{2})}+O(D^{-\infty}).
\]
The oscillation is evaluated in closed form, with no cancellation estimate. Since
$C^{-u}D^{-u/2}=(C^2D)^{-u/2}$, collecting the $u$-contour integral gives the $m$-th term
$e(D)(2\pi iD)^{-(m+1+a/2)}A^+_m(C^2D)$ with $A^+_m$ exactly \eqref{eq:A6}: the coefficients $b_m(u)$,
carrying the $(2-t)^{u/2}$ Taylor data, reproduce Altu\u{g}'s $c^\pm_m$. (One identifies the two Taylor
systems term by term: $b_m(u)$ and $c^\pm_m(u/2)$ are both the $m$-th coefficient in the $(2-t)^{u/2}$
expansion of the endpoint profile $\varphi(t)(2-t)^{u/2}\chi_+(1-t)$, so they coincide.)

\emph{Remainder.} The Taylor remainder satisfies $|R_M(u,t)|\le C_M(u)\,t^{M+1}$ with
$C_M(u)=\sup_{[0,\delta]}|\partial_t^{M+1}\psi_u|/(M+1)!$ of at most polynomial growth in $\Im u$
($\psi_u$ smooth). Its contribution is again a transform of the form \eqref{eq:A6} with amplitude
$R_M$; taking its $u$-contour at $\Re u=\tau_1$ (any $\tau_1>0$), Lemma~\ref{lem:mag} bounds it by
$\mathcal B_{M,\tau_1}(C^2D)^{-\tau_1/2}D^{-(M+1+a/2)}$ with $\mathcal B_{M,\tau_1}<\infty$
independent of $C,D$. The endpoint $x=-1$ gives the $e(-D)$, $A^-_m$ terms identically.
\end{proof}

\begin{remark}[on the exponents]
The contour height $\tau_1>0$ is free, and the remainder decays like $(C^2D)^{-\tau_1/2}$ because the
factor $(C^2D)^{-u/2}$ is evaluated at $\Re u=\tau_1$; Altu\u{g}'s stated $(C^2D)^{-\tau_1}$ is the same
statement in his normalization of the Mellin variable ($u$ versus $u/2$). What the application uses is
only that, for arbitrary $A>0$, the remainder is $O_A\big((C^2D)^{-A}D^{-(M+1+a/2)}\big)$ uniformly in
$C,D$ --- immediate from the freedom in $\tau_1$ together with Lemma~\ref{lem:mag}.
\end{remark}

\begin{proposition}[Part~III Prop.~5.2]\label{prop:52}
$|J_{l,f}(\xi,\nu,X)|\ll X^2\nu^{-M}(lf^2)^{-N}\big((lf^2/\sqrt X)^{N-M+3}+\xi^M\big)$ with implied
constant independent of $(\xi,\nu,l,f,X)$.
\end{proposition}

\begin{proof}
Write $J_{l,f}(\xi,\nu,X)=\int G(y/X)\,y\,\Psi(y)\,e(-y\nu/2lf^2)\,dy$ with $\Psi(y)=I_{l,f}(\xi,y)$.
Since $G$ is compactly supported, integrating by parts $M$ times against the phase gives the exact
identity $J=(lf^2/(-\pi i\nu))^M\int\partial_y^M(G(y/X)\,y\,\Psi(y))\,e(-y\nu/2lf^2)\,dy$, hence
\[
  |J_{l,f}(\xi,\nu,X)|\le\Big(\tfrac{lf^2}{\pi\nu}\Big)^{M}
    \big\|\partial_y^M\!\big(G(y/X)\,y\,\Psi(y)\big)\big\|_{L^1(y)} .
\]
By the Leibniz rule
$\partial_y^M(G(y/X)\,y\,\Psi)=\sum_{i+j+k=M}\binom{M}{i,j,k}X^{-i}G^{(i)}(y/X)\,(y)^{(j)}\,
\partial_y^k\Psi$, where $(y)^{(j)}=y,1,0$ for $j=0,1,\ge2$. The cutoff $G^{(i)}(y/X)$ is supported on
$y\in[\tfrac X4,\tfrac{5X}4]$ (length $\asymp X$) and contributes $\|G^{(i)}\|_\infty$, and the factor
$(y)^{(0)}=y\asymp X$ supplies a second power of $X$, giving the $X^2$. Each $y$-derivative of
$\Psi(y)=\int\theta_\infty(x)[F+H](\arg)\,e(-x\xi\sqrt{4y}/4lf^2)\,dx$ falls on either the phase ---
producing a factor $O(\xi/lf^2\sqrt y)$, the source of the $\xi^M$ branch --- or on the profile through
$\arg=lf^2/(\sqrt{4y}\sqrt{1-x^2})$, producing a factor $O(lf^2/y)$, the source of the $(lf^2/\sqrt
X)$ branch. Each resulting $x$-integral is a transform of the form \eqref{eq:A6} (with $C\asymp\xi\sqrt
y/lf^2$ and $D$ the local scale), bounded by Lemma~\ref{lem:mag}, or at the sharp order by the
expansion of Theorem~\ref{thm:A14}. Collecting the $X$-, $lf^2$-, and $\xi$-powers reproduces the two
branches $(lf^2/\sqrt X)^{N-M+3}$ and $\xi^M$; the implied constant is a finite sum of the integrals
$B_{\tau,m,a}(\Phi)$ (Lemma~\ref{lem:mag}) and the norms $\|G^{(i)}\|_\infty$, none depending on
$(\xi,\nu,l,f,X)$.
\end{proof}

Every uniform constant in \S\ref{sec:gauge} is one application of Lemma~\ref{lem:mag}; the content of
Altu\u{g}'s Appendix~A, from the exact-gauge standpoint, is that single magnitude estimate. Where
Altu\u{g}'s route controls the transform through the asymptotic expansion together with an oscillation
bound on the remainder, the exact gauge replaces that oscillation bound by the triangle inequality in
the coordinate that makes the integrand real: the whole $(\xi,\nu,l,f,X)$-uniformity then rests on the
$x$-independence of the single line integral $B_{\tau,m,a}(\Phi)$, with no cancellation estimate at any
step.

\section{The oscillation is the cutoff's comb}\label{sec:comb}

Why did Prop.~\ref{prop:52} look hard? The oscillation the estimate must survive is deterministic.
Scanning $|J_{l,f}(\xi,\nu,X)|$ in the transform frequency $\nu$ over the cutoff support
$y\in[\tfrac X4,\tfrac{5X}4]$ reveals a comb of near-zeros with spacing, to measurement,
$\Delta\nu=4lf^2/(\kappa X)$ with $\kappa\approx0.94$ the support fraction, matching the two-edge
beat of $G$: the dispersion law $\Delta\nu\propto(lf^2)^{1.00}(X/8)^{-0.98}(1+\xi)^{0.02}$ fits at
$R^2\approx1$ (Appendix~\ref{app:repro}) --- the interference of the two edges of the fixed cutoff, switched on once
the $\xi$-phase drives the $x$-integral into edge dominance. At $\xi=0$ the $x$-integrand is unimodal
(prominence ratio $P_2/P_1=0$) and the comb is absent. The apparent difficulty is a chart feature of
the fixed cutoff---exactly what the exact gauge of \S\ref{sec:gauge} removes.

\section{The productivity boundary}\label{sec:sarnak}

The symmetric-power kernel is modulated by $U_r(\cos\theta)=\sum_{j=0}^r e^{i(r-2j)\theta}$, carrying
a zero-frequency component iff $r$ is even. At $\xi=0$ the $x$-integral reads this component (the
detecting residue); at $\xi\neq0$ it is the removable comb of \S\ref{sec:comb}, bounded by
Prop.~\ref{prop:52} (with $|U_r|\le r+1$) and hence $o(X)$. The two regions are disjoint in $\xi$:
the detecting residue sits at $\xi=0$, where the comb is absent; the comb sits at $\xi\neq0$, where it
is removable. So the residue isolates for every $r$. On the dual integral (level $p^k=25$) the
zero-frequency value against the $\xi\neq0$ floor is
\[
  r=1:\ 0\ (\text{exact}),\quad r=2:\ 55.9,\quad r=3:\ 1.4\times10^{-16},\quad r=4:\ 11.4,
\]
nonzero iff $r$ even, clear of the floor by an order of magnitude. The degradation with $r$ is the
growth of the character tail mass ($0.042,0.092,0.135,0.156,0.256$ for $r=0,\dots,4$, fit
$\text{ceiling}\times\sqrt{\#\text{harmonics}}$ at $R^2=0.977$)---gradual attrition, not a wall.
Sarnak's productivity boundary \cite[fn.~5]{AltIII} records the growth of a removable floor, not a
loss of the residue.

\section{Toward functoriality: the product law and the transfer factor}\label{sec:gue}

The three preceding sections close the analytic core --- the uniformity (\S\ref{sec:gauge}), the
emergent clock (\S\ref{sec:comb}), and the productivity boundary. The remainder of the paper develops
the two mechanisms toward symmetric-power functoriality, beginning with the product law, which supplies
the transfer factor.

The derivative statistic that random-matrix theory reads at an eigenvalue is the reopening rate an
arithmetic fiber reads at a zero, by \eqref{eq:prod}: both are $\prod_{r_j\neq\mu}|\mu-r_j|$. This is
an identity, not a coincidence, and it is the mechanism the ensemble leaves implicit. The consequence
here is structural: the same distance product is a \emph{Weyl denominator}, and Weyl denominators are
the archimedean transfer factors of \S\ref{sec:fl}. Self-interference is one object appearing as a
zero statistic, a reopening rate, and a transfer factor.

\begin{remark}[a distinct post-closure mechanism, deferred]
The product law fixes the derivative statistic \eqref{eq:prod} but not the \emph{configuration} that
feeds it, and that configuration is governed by a separate dynamics we only flag here. At a focal
cancellation the ordinary phasor residue vanishes, yet the fiber carries a nonzero eigen response that
persists over roughly one native harmonic cell; when closures fall within a cell their response
intervals overlap and deform the local spacing deterministically. The algebraic backbone --- a
\emph{cluster product law} generalizing \eqref{eq:prod} to any finite cluster --- is proven
(\texttt{ReverbResidue.cluster\_product\_law}, with pair suppression and the reopening residue), and
the dynamics is measured to hold uniformly across pairs, triples, quads, and clusters carrying
non-uniformly spaced neighbours: it is not a pairwise artifact, the case a triple would expose. This is
a deterministic candidate for the configuration law that populates \eqref{eq:prod} --- like the
unitary ensembles, but with a reason. We make no limiting-ensemble claim here (finite-$N$
characteristic-polynomial geometry is naturally unitary-circular, while the unfolded zero statistics
live in the limiting language); formulating the response kernel and deriving the unfolded $n$-level
correlations is the subject of a separate paper.
\end{remark}

\section{The carrier/fiber rescaling: a pole becomes a closed cell}\label{sec:carrier}

Place the integers on a carrier scaled by $\pi/3$, i.e. attach the cell phase $e^{in\pi/3}$ (the
$\mu_6$/$\zeta_6$ step). The pole of $\zeta$ at $s=1$ lives in the divergence of $\sum 1/n$; on the
carrier it closes,
\[
  \sum_{n\ge1}\frac{e^{in\pi/3}}{n}=-\log\!\big(1-e^{i\pi/3}\big)=\frac{i\pi}{3},
\]
a finite $\mu_6$ value: the divergence has become a closed-cell reading. For a Dirichlet character
$\chi$ of conductor $q$, the appropriate fiber scale (trivial $\to\pi/6$, $\chi_3\to\pi/3$,
$\chi_4\to\pi/2$) makes carrier and fiber compatible harmonics; the principal character, carrying the
pole, is handled by the eta regulator,
\[
  \sum_{n\ge1}\frac{(-1)^{n-1}e^{in\pi/3}}{n}=\log\!\big(1+e^{i\pi/3}\big)=\tfrac12\ln3+\tfrac{i\pi}{6},
\]
landing on the trivial character's own cell ($\pi/6$, $\mu_{12}$), pole-free, with the partial-sum
phase converging (no drift). In this frame the $S(t)$ oscillation of the readout is removed---the
value is a fixed harmonic, not a wandering phase---so what was a residue to be identified is a
zero-frequency cell value, read as the DC-residue $(s-c)F'/F\to r$ of a logarithmic derivative. This
is the reading used in \S\ref{sec:detect}.

\section{Symmetric-power transfer: arithmetic transparency}\label{sec:fl}

For $\Sym^r$ transfer $\mathrm{GL}(2)\to\mathrm{GL}(r+1)$ ($r=2m$ even), the transfer sends an
elliptic class of Satake angle $\theta$ to the class of angles $\{2k\theta:|k|\le m\}\cup\{0\}$
---i.e. $\Sym^r$ of the Frobenius block, with the fixed eigenvalue $1=\alpha\beta$ the det-one of
the block. Two features make the archimedean transfer factor explicit and the arithmetic transparent.

\emph{Arithmetic transparency.} The transferred characteristic polynomial is reducible ($x-1$
divides it), and each quadratic factor $x^2-2\cos(2k\theta)x+1$ has discriminant $-4\sin^2(2k\theta)$.
The clock identity $\sin(2k\theta)=\sin\theta\,U_{2k-1}(\cos\theta)$ (with $\sin^2\theta=|D|/4n$,
$D=t^2-4n$) gives
\[
  \frac{\operatorname{disc}_k}{D}=\frac{U_{2k-1}(\cos\theta)^2}{4n}\in\mathbb Q^{\times2},
\]
a rational square---so every factor lies in the \emph{same} field $\mathbb Q(\sqrt D)$ as the
original class, the arithmetic $L$-value is identical on both sides, and it cancels in the transfer
factor. (Because symmetric powers transfer to \emph{reducible} classes, this avoids the order-zeta
conjecture of Deng--Espinosa, which concerns irreducible cubic orders.)

\emph{Transfer factor as a distance product.} What remains is the ratio of Weyl denominators
---distance products by \eqref{eq:prod}. For $r=2$,
\[
  \frac{\Delta_3}{\Delta_2}
  =\frac{|1-e^{2i\theta}|^2\,|e^{2i\theta}-e^{-2i\theta}|}{|e^{i\theta}-e^{-i\theta}|}
  =4|\sin\theta||\sin2\theta|=8\sin^2\theta|\cos\theta|=\frac{|D|\,|t|}{n^{3/2}},
\]
a monomial in the class data; its winding phase, gauged by the $\pi/3$ ($\mu_6$) step, is a real
harmonic (\S\ref{sec:carrier}). Thus $\Delta(\theta)=U_r(\cos\theta)\cdot(\Delta_{r+1}/\Delta_2)
\cdot(\text{radial }\sqrt p\text{-scaling})$, the Weyl character and distance product being product-law
objects and the phase an exact-gauge object.

\section{Detection: the self-twist identity closes}\label{sec:detect}

We isolate one computable instance of the detection step and close it on a genuine transfer. For a
holomorphic newform $f$ with $a_p$ its Frobenius trace and $\lambda(p)=a_p/p^{(k-1)/2}$, the
symmetric square carries the middle eigenvalue $\alpha\beta=1$; the plain census
$\sum_p(\lambda(p)^2-1)(\log p)/p$ has mean zero and is blind. The correct object is the
\emph{self-twist} Rankin--Selberg $L$-function $L(s,f\times f\otimes\chi_K)$ for the quadratic
character $\chi_K$ of an imaginary quadratic field $K$: it has a pole at $s=1$ iff $f$ has complex
multiplication by $K$ (dihedral), the classical criterion of Gelbart--Jacquet \cite{GJ} and
Rankin--Selberg.
Its pole order is the transferred count, read as the zero-frequency residue of the logarithmic
derivative:
\[
  \text{(LEFT, analytic)}\qquad
  A_{\mathrm{tw}}(X)=\sum_{p\le X}\lambda(p)^2\chi_K(p)\frac{\log p}{p},\qquad
  \frac{d\,A_{\mathrm{tw}}}{d\log X}\ \longrightarrow\ \ord_{s=1}L(s,f\times f\otimes\chi_K).
\]
The geometric side is the self-twist relation on Frobenius, $f\otimes\chi_K\cong f$, equivalently
$a_p=0$ at every inert prime:
\[
  \text{(RIGHT, geometric)}\qquad
  \#\{p\le X\ \text{inert}:\ a_p=0\}\big/\#\{p\le X\ \text{inert}\}\ \longrightarrow\ 1_{\{f\ \text{has CM by }K\}}.
\]

We evaluate both on two weight-two forms with $K=\mathbb Q(\sqrt{-3})$ (the Eisenstein field, whose
$\chi_K(p)=+1,-1$ for $p\equiv1,2\pmod3$): the CM curve $y^2=x^3+1$ and the non-CM curve
$y^2=x^3-16x+16$ (conductor $37$), with $a_p$ from point counting.

\begin{center}
\begin{tabular}{lccc}
\toprule
 & LEFT: slope $dA_{\mathrm{tw}}/d\log X$ & RIGHT: inert $a_p=0$ fraction & count \\
\midrule
CM $\sqrt{-3}$ ($x^3+1$) & $0.997$ & $725/726=0.9986$ & $1$ \\
non-CM ($x^3-16x+16$) & $0.037$ & $6/726=0.0083$ & $0$ \\
\bottomrule
\end{tabular}
\end{center}

For the CM form the analytic slope is $0.97$--$0.99$ across every interval $[1000,4000]$,
$[4000,10000]$, $[10000,20000]$ (pole order $1$), and the geometric self-twist holds at every good
inert prime (the single exception $725/726$ is $p=2$, the bad-reduction prime of conductor
$36=2^2\cdot3^2$). For the non-CM form both sides are $0$. Thus
\[
  \boxed{\ \ord_{s=1}L(s,f\times f\otimes\chi_K)\ =\ 1_{\{f\otimes\chi_K\cong f\}}\ =\ \text{transferred count}\ }
\]
closes---$1$ on a genuine transfer, $0$ on a genuine non-transfer---read through the logarithmic-
derivative DC-residue on the S(t)-free carrier of \S\ref{sec:carrier}, with the emergent-clock
lesson explaining why the plain census cancels (inert primes carry $\lambda^2=0$, so the self-twist
census has no cancelling floor).

\begin{remark}[what this is]
The mathematics detected here---that the symmetric-square/self-twist pole marks CM forms---is
classical (Gelbart--Jacquet, Rankin--Selberg). What the section contributes is the closure of the
detection \emph{identity} through the two mechanisms: the analytic count as a product-law/log-
derivative residue, the geometric count as the self-twist relation, matched on a genuine transfer.
\end{remark}

\subsection{A non-abelian instance: which Frobenius data transfers}\label{sec:a4}

The self-twist detection of \S\ref{sec:detect} runs on a dihedral (abelian, $C_2$) form, where the
Galois image is small and the arithmetic reduces to a single quadratic character. For a genuinely
non-abelian form the Frobenius datum splits into two independent bits---one a character cannot
see---and the transfer treats them differently. We illustrate this on the exotic tetrahedral
(``$A_4$'') weight-one newform $h$ of level $124$, whose $L$-coefficients are a class function of
$\mathrm{Frob}_p$ in the binary tetrahedral group $2T=\mathrm{SL}(2,3)$, given by the closed formula
$a_p=\zeta_{12}^{k(p)}$ with
\[
  k(p)=\tfrac12\big(\chi_k(p)+4\varepsilon(p)\big)+\big(0\ \text{if }(3/p)=+1,\ \text{else }6\big)
  \pmod{12},
\]
$a_p=0$ when $\varepsilon(p)=0$. Here the exponent carries two independent Galois bits:
\[
  \text{\emph{value bit}}\ \varepsilon(p)\in\{0,\pm1\}\ \text{(the cubic residue, the abelian
  $A_4^{\mathrm{ab}}=C_3$ datum)},\qquad
  \text{\emph{sign bit}}\ (3/p)\ \text{(the $\mathrm{SL}(2,3)\!\to\!A_4$ central $\pm1$,}
\]
the double-cover datum read from the order of $\mathrm{Frob}_p$; the sign bit resolves $k$ versus
$k+6$ and is invisible to any abelian character.

The symmetric square sends $a_p\mapsto a_p^2-\chi(p)$, and $a_p^2=\zeta_{12}^{2k}$ is unchanged by
$k\mapsto k+6$. Hence:

\begin{center}
\begin{tabular}{lc}
\toprule
 & Sym$^2$ transfer to GL(3) \\
\midrule
sign bit $(3/p)$ (double cover, $\mathrm{SL}(2,3)$) & \emph{forgotten} (trace invariant under $k\!\leftrightarrow\!k{+}6$) \\
value bit $\varepsilon$ (cubic residue, projective $A_4$) & \emph{retained}: $\varepsilon=\pm1\Rightarrow\pm\sqrt3\,i$, $\varepsilon=0\Rightarrow\pm1$ \\
\bottomrule
\end{tabular}
\end{center}

So the transfer factors through the projective image $A_4$: the value bit passes to the GL(3) Satake
datum, the sign bit is the GL(2)-specific double-cover fiber the transfer collapses. The order trick
---sign from the element order, value from the residue symbol---thus reads \emph{exactly which part
of the Frobenius datum transfers}, and the split is non-trivial precisely because $A_4$ is
non-abelian (for the dihedral case of \S\ref{sec:detect} the two bits coincide). This identifies the
transferred datum; the automorphy of the resulting GL(3) object (the three-dimensional standard
representation of $A_4$) is the classical Artin/Langlands--Tunnell case.

\subsection{The detected count is a Mordell--Weil rank difference}\label{sec:rankdiff}

The transferred object of \S\ref{sec:a4} is, up to the nebentype twist, the three-dimensional
adjoint $\mathrm{Ad}_g=\mathrm{End}^0(V_g)\cong\Sym^2 V_g\otimes(\det V_g)^{-1}$ of the $A_4$ form
$g$. At this exotic locus the zero-frequency count read on the analytic side coincides with a
genuine arithmetic rank, and the coincidence is an exact identity of group theory.

The quartic field $M$ cut out by $g$ carries the permutation action of $A_4$ on four points, whose
character is $\chi_{\mathrm{perm}}=(4,0,1,1)$ on the classes $(e,\ (2,2),\ (3)_a,\ (3)_b)$. Against the
character table, $\langle\chi_{\mathrm{perm}},\mathbf 1\rangle=1$ and
$\langle\chi_{\mathrm{perm}},\chi_{\mathrm{std}}\rangle=1$, so
\[
  \chi_{\mathrm{perm}}=\mathbf 1\ \oplus\ \mathrm{Ad}_g,\qquad \mathrm{Ad}_g=\chi_{\mathrm{std}}\
  \text{(the $3$-dimensional irreducible)} .
\]
Functoriality of Mordell--Weil under this decomposition gives
$E(M)\otimes\mathbb Q = (E(\mathbb Q)\otimes\mathbb Q)\oplus E(\,\cdot\,)_{\mathrm{Ad}_g}$ and hence the
\emph{exact} Artin/permutation identity
\begin{equation}\label{eq:rankdiff}
  r(E,\mathrm{Ad}_g)\ =\ \operatorname{rank}E(M)-\operatorname{rank}E(\mathbb Q),
\end{equation}
proven as group theory (no analytic or $p$-adic input). The left side is the order of vanishing at
$s=1$ of $L(E,\mathrm{Ad}_g,s)$---the zero-frequency residue of the logarithmic derivative, the same
statistic \S\ref{sec:detect} read as the transferred count. The right side is a difference of
classical Mordell--Weil ranks, the source living ``up the tower'' at $M$ and traced down to
$\mathbb Q$. Thus the count is a rank difference: one statistic, read once on the automorphic side
and once on the arithmetic side.

On the two Darmon--Lauder--Rotger exotic $A_4$ curves this is pinned on both sides:
\begin{center}
\begin{tabular}{lccc}
\toprule
 & $\operatorname{rank}E(\mathbb Q)$ & $\operatorname{rank}E(M)$ & $r(E,\mathrm{Ad}_g)=\operatorname{rank}E(M)-\operatorname{rank}E(\mathbb Q)$ \\
\midrule
$26b$ ($y^2+xy+y=x^3-x^2-3x+3$) & $0$ & $2$ & $2$ \\
$52b$ ($y^2=x^3+x-10$) & $0$ & $2$ & $2$ \\
\bottomrule
\end{tabular}
\end{center}
with $M=\mathbb Q(w)$, $w^4+7w^2-2w+14=0$ the tetrahedral field of \S\ref{sec:a4}, rank-two
generators of $E(M)$ exhibited by descent. This is the ``one mechanism, two locations'' statement
made theorem-anchored: at the exotic locus both sides of \eqref{eq:rankdiff} are determined, the
identity between them is proven group theory, and the shared statistic is the DC-residue that reads
the functorial pole in \S\ref{sec:detect} and the analytic rank in the Beyond-Endoscopy census
alike.

\begin{remark}[tiering]
The rank identity \eqref{eq:rankdiff} is proven (Artin formalism). The automorphy of $\mathrm{Ad}_g$
on GL(3)---which makes $L(E,\mathrm{Ad}_g,s)$ an honest automorphic $L$-function so that its pole order
is a spectral quantity---is the classical Langlands--Tunnell case for $A_4$. The $p$-adic
\emph{regulator} side (the explicit point coordinates realizing the rank) is the Darmon--Lauder--Rotger
Elliptic Stark conjecture, verified to $20$ digits on $26b$/$52b$ but with the points found by
descent first; we use only the proven rank identity and the classical automorphy here.
\end{remark}

\section{A candidate automorphic object from the primary construct}\label{sec:construct}

The transfer of \S\S\ref{sec:fl}--\ref{sec:rankdiff} requires an automorphic representation on
$\mathrm{GL}(r+1)$ whose existence realizes the functorial lift. In the primary (three-dimensional)
representation we exhibit a concrete \emph{candidate object} for it: the \emph{carrier} --- a
double-sided helix of Frobenius similitudes --- carrying a \emph{fiber} (the phasor readout), together
with its \emph{warp} (the twist family of \S\ref{sec:fl}). The classical $L$-function is the shadow of
this candidate under the projection that sends the exponential state space $e^{y}$ down by a
M\"obius/Cayley map and then $\log$ to the ordinate, at the midpoint and without drift. What makes it
a candidate is that the three properties a converse theorem asks of an $L$-function --- functional
equation, analytic continuation, and boundedness/completeness --- are each a \emph{proven} statement
about it, and each is already formalized.

\paragraph{Functional equation: the reflection axis.} The anti-helix is the complex \emph{conjugate}
of the helix, not its mirror: the carrier's transverse block is $\operatorname{diag}(z,\bar z)$ with
$z$ the Frobenius spin. It is literally the de Branges conjugate-pair block
(\texttt{frobeniusBlock\_eq\_conjPairBlock}), with $z\bar z=|z|^2=1$ (\texttt{spin\_mul\_conj}), so
$\det=1$ and it is unitary (\texttt{frobeniusBlock\_det\_one}, \texttt{frobeniusBlock\_unitary}). The
det-one condition is not a normalization but the \emph{reality locus}:
$\det=1\iff|z|=1\iff$ the helix and anti-helix \emph{balance} $\|E(z)\|=\|E^{*}(z)\|\iff z$ real
(\texttt{frobeniusBlock\_deBranges\_reality\_bridge}). The dualizing reflection $s\mapsto1-s$ fixes
\emph{exactly} the critical line, $\Re(1-s)=\Re s\iff\Re s=\tfrac12$ (\texttt{reflection\_fixes\_iff}),
and this abscissa is not posited but \emph{forced} by the arclength area law --- the carrier admits a
scale-balanced fiber iff $\sigma=\tfrac12$ (\texttt{scaleBalanced\_iff}), the reciprocation of the
fiber amplitude $n^{-\sigma}$ against the area-law radius $\sqrt n$. This is the functional equation
realized geometrically: the conjugate leg closing on the reflection-fixed axis. It survives
\emph{every} unit-modulus warp $z\mapsto Az$ ($|A|=1$), $z\bar z\mapsto|A|^2=1$
(\texttt{warpedBlock\_det\_one\_of\_warp}), so the whole twist family carries it, structurally and
without estimate. (A complex twist is \emph{self}-dual only when it equals its own conjugate;
otherwise the equation pairs the sideband $\chi_j$ with $\chi_{-j}=\bar\chi_j$, as the conjugate leg
demands.)

\paragraph{The crossing geometry and the two fiber topologies.} A zero is a phasor \emph{focal
cancellation}: the phasors, generated continuously as the carrier grows, align and cancel at a height
fixed by the carrier scale, the order of the integers/primes, and the conductor. Because the same data
drive both conjugate legs, the two legs cancel at the \emph{same} height, the determinant-one link
holding between them. Read as its own oscillator the fiber accumulates phase along each leg; the
crossing spacing is a half-turn on each side of the origin and a full turn thereafter. This reduces to
the functional equation making the central fiber an extremum --- $Z$ real and even, $Z'(0)=0$
(\texttt{collapseWave\_even}, \texttt{hinge\_turning\_point}) --- together with the argument principle,
so the half-turn offset and full-turn step are provable; only the \emph{rigidity} of the spacing in $t$
stays RH-adjacent. The radial envelope of the winding is the Archimedean spiral, whose arclength area
law $r_n\sim\sqrt n$ forces the abscissa to $\tfrac12$ (\texttt{scaleBalanced\_iff}).

Two fiber topologies occur. An \emph{open} helix --- the Dirichlet family --- has its fibers begin at
the origin and run up both conjugate legs to infinity. A \emph{clipped} helix --- a self-dual
$L$-function such as that of an elliptic curve --- has its strand bounded by the conductor: the live
phasor count scales as $\sqrt N$, the analytic conductor (measured exponent $0.5000$ across seven
decades, Appendix~\ref{app:repro}), the same $\sqrt{\,\cdot\,}$ law that fixes the abscissa; and its
two legs \emph{meet} at the central point --- the fixed point of the functional-equation involution
(UNIT/2, \texttt{affine\_reflection\_fixed\_iff}), the ordinate origin $t=0$ (equivalently
$s=\tfrac12$ analytic, $s=1$ arithmetic), not the spiral base $N=0$. There the two strands carry the
determinant-one conjugate weights $r^{s-1}\!\cdot r^{1-s}=1$ (\texttt{strand\_weights\_det\_one}), and
the root number $\varepsilon$ \emph{is} the weld sign: $\varepsilon=-1$ cancels the central term and
forces a central zero (\texttt{weld\_kills\_each\_phasor}), $\varepsilon=+1$ doubles it
(\texttt{weld\_doubles\_each\_phasor}). The same double-ended weld reads, in one factorization
$(s-c)^r G$, the root number, the rank $r$ (the DC-residue, \texttt{rank\_is\_dc\_residue}), and the
leading datum $L^{(r)}/r!$ --- matched to tabulated values on ranks $0$–$3$.

\paragraph{Continuation: the readout, continued past the strip.} The fiber's value is the phasor
representation $L(\chi,s)=\sum_n\chi(n)\,n^{-\sigma}\,\mathrm{mellinSpin}(y,n)$
(\texttt{LSeries\_phasor\_representation}), and it is genuinely \emph{continued}, not merely
re-indexed. For a non-principal character the buckets cancel over each period, so the partial sums
$\sum_{n<N}\chi(n)$ are bounded (\texttt{character\_partialSum\_norm\_le}); Abel summation against the
bounded-variation weight $n^{-s}$ then carries the readout past the $\Re=1$ absolute-convergence wall,
and the Abel-summed tail --- which agrees with the $L$-function on $\Re s>1$ and is analytic on the
connected strip $\Re s>0$ --- is forced by the identity theorem to converge to $L(\chi,s)$ on the
\emph{whole} strip $\Re s>0$ (\texttt{dirichlet\_strip\_tendsto\_LFunction}). The principal/$\zeta$
case runs identically through the eta regulator on the punctured strip $\{\Re s>0\}\setminus\{1\}$,
correctly excising the pole (\texttt{eta\_strip\_tendsto}). This is a genuine analytic continuation,
machine-checked, with no remainder: the finite harmonic cells (the $\mu_6$ buckets) cancel exactly.
The completed function $\Lambda(s)=\gamma(s)L(s)$ and the \emph{global} functional equation are then
the standard Tate \cite{Tate}/Godement--Jacquet \cite{GJa} assembly of the per-place det-one links by
Poisson summation on $\mathbb A/\mathbb Q$.

\paragraph{Boundedness and completeness.} Two inputs, both discharged in the construct.
\emph{Bounded in vertical strips}: the Abel tail of the previous paragraph is locally uniformly
summable on $\{\Re s>0\}$ (\texttt{summableLocallyUniformlyOn\_tail}), so the continued readout is
bounded on every compact of the strip --- the Phragm\'en--Lindel\"of conclusion, obtained natively from
the bounded (finite-cell) partial sums rather than from a convexity estimate.
\emph{Complete}: the object has no spectral mass unaccounted for. The event space of the
three-dimensional object is the height ray, which is the weld; the strip is a one-dimensional
projection device with no primary counterpart, so there is nowhere off-weld for a zero event to hide,
and every zero event is weld-fixed and is a source --- unconditionally, for \emph{every} fiber:
\[
  \texttt{threeD\_exhaustive}\ (E:\mathbb C\to\mathbb C):\quad
  \forall\,y\in\mathbb R,\ E(y)=0\ \Rightarrow\ \mathrm{IsSource}\,E\,(y).
\]

\paragraph{The candidate, assembled.} On the primary construct the object exists; its functional
equation is the reflection-fixed axis surviving every warp, its continuation is the readout carried to
$L(\chi,\cdot)$ across the whole strip, its boundedness is the local-uniform tail, and its
completeness is exhaustion --- each a proven, formalized statement. The carrier, fiber, and warp
\emph{are} the candidate automorphic object, and the converse-theorem inputs are discharged in the
primary representation. These inputs --- the warp-invariant det-one link and the reflection axis
$\Re s=\tfrac12$ (angle), the phasor readout continued past the strip (continuation), three-dimensional
exhaustion (height/completeness), the midpoint landing on $\Re s=\tfrac12$ with the admissible locus
forced to $\{\tfrac12\}$, and the loss-ledger reconstruction bijection (recovery) --- are bundled and
machine-checked in a single theorem (\texttt{AutomorphicCandidate.candidate\_wellformed}), with the
standard axiom footprint. This is the object we use to realize the functorial lift of
\S\S\ref{sec:fl}--\ref{sec:rankdiff}; the bridge from it to the classical statement that $\Sym^r\pi$ is
automorphic on $\mathrm{GL}(r+1)$ is the converse theorem of \S\ref{sec:converse}, whose hypotheses are
exactly these proven properties.

\paragraph{The local identification.} The properties above are properties of an $L$-function; to name
that $L$-function as $L(s,\Sym^r\pi)$ we record, once and explicitly, the equality of local data. This
is the equality the converse theorem consumes, and it is fixed by the construction rather than by any
assumption on the target.

\begin{proposition}[Global identification of the candidate with $\Sym^r\pi$]\label{prop:localid}
Let $\pi$ be a cuspidal automorphic representation of $\mathrm{GL}(2)/\mathbb Q$, and for each place $v$
let $\varphi_{\pi,v}\colon W'_{\mathbb Q_v}\to\mathrm{GL}_2(\mathbb C)$ be its local Langlands parameter
(a theorem of local Langlands for $\mathrm{GL}(2)$). At \emph{every} place $v$ the standard local factor
and local root number of the candidate object of \S\ref{sec:construct} are those of the $(r{+}1)$-dimen\-sional
parameter $\Sym^r\varphi_{\pi,v}$:
\begin{equation}\label{eq:localid}
  L_v(s,\mathrm{cand})=L\!\left(s,\Sym^r\varphi_{\pi,v}\right),\qquad
  \varepsilon_v(s,\mathrm{cand},\psi_v)=\varepsilon\!\left(s,\Sym^r\varphi_{\pi,v},\psi_v\right).
\end{equation}
At an unramified $v$, $\varphi_{\pi,v}=A_{\pi,v}=\operatorname{diag}(\alpha_v,\beta_v)$ ($\alpha_v\beta_v
=1$) and \eqref{eq:localid} reads $L_v=\det(1-\Sym^r(A_{\pi,v})q_v^{-s})^{-1}=\prod_{j=0}^{r}(1-\alpha_v^{
\,r-j}\beta_v^{\,j}q_v^{-s})^{-1}$; at $v=\infty$ the parameter $\Sym^r\varphi_{\pi,\infty}$ gives
$L_\infty=\prod_i\Gamma_{\mathbb C}(s+\mu_i)$ (one $\Gamma_{\mathbb R}(s+\delta)$ adjoined for even $r$;
for $\pi$ of weight $k$, $\mu_i=(2i{-}1)\tfrac{k-1}2$). Consequently the \emph{full} completed
$L$-function, the global root number $\varepsilon=\prod_v\varepsilon_v$, and the conductor
$N=\prod_v q_v^{\,a(\Sym^r\varphi_{\pi,v})}$ of the candidate are exactly those of $\Sym^r\pi$ --- the
entire global package, at all places, ramified included. No automorphy of $\Sym^r\pi$ is used:
$\varphi_{\pi,v}$ is furnished by local Langlands for $\mathrm{GL}(2)$, and $\Sym^r$ is a functor on
parameters.
\end{proposition}

\begin{proof}
The candidate carrier is assembled from the Frobenius similitudes of $\pi$, so its transverse datum at a
place $v$ is exactly $\varphi_{\pi,v}$ (at unramified $v$ the block $A_{\pi,v}=\operatorname{diag}(z_v,
\bar z_v)$, $z_v\bar z_v=1$, \texttt{frobeniusBlock\_eq\_conjPairBlock}). The fiber over $v$ is, by the
transfer of \S\ref{sec:fl}, the $r$-th symmetric power of that datum, $\Sym^r\varphi_{\pi,v}$; its
standard local $L$- and $\varepsilon$-factors are the local Langlands/Deligne factors of the
Weil--Deligne representation $\Sym^r\varphi_{\pi,v}$ \cite{Tate}, which is \eqref{eq:localid}. At
unramified $v$ this is the reciprocal characteristic polynomial of
$\operatorname{diag}(\alpha_v^{r},\dots,\beta_v^{r})$, $\det=(\alpha_v\beta_v)^{r(r+1)/2}=1$; at $\infty$
it is the $\Gamma$-datum of $\Sym^r\varphi_{\pi,\infty}$ read by the midpoint projection
(\texttt{midpoint\_projects\_to\_half}). The three global invariants are the corresponding products over
$v$ (finite for $\varepsilon$ and $N$, since $\varphi_{\pi,v}$ is unramified for almost all $v$), and
these coincide, place by place, with the standard data of $\Sym^r\pi$. Nowhere is $\Sym^r\pi$ assumed
automorphic: $\varphi_{\pi,v}$ is the local parameter of the given $\mathrm{GL}(2)$ representation $\pi$,
and $\Sym^r\varphi_{\pi,v}$ is a definite $(r{+}1)$-dimensional local parameter.
\end{proof}

Proposition~\ref{prop:localid} upgrades the \emph{identification} to all places: the candidate's
$L$-function, root number, and conductor \emph{are} those of $\Sym^r\pi$, globally, with the exact
$\varepsilon$/conductor package and no automorphy assumed. It does not upgrade the \emph{niceness}: that
the completed $L(s,\Sym^r\pi\times\sigma)$ satisfies its functional equation is the analytic input the
candidate supplies on the standard side (\S\ref{sec:construct}) and that, for the twisted convolutions,
remains the converse-theorem hypothesis of \S\ref{sec:converse} --- unconditional for $r\le4$
(Langlands--Shahidi), the single global assembly outstanding for $r\ge5$. Identification is now global;
niceness is where the frontier stays.

\paragraph{Recovery: the loss-ledger round trip.} The projection to the shadow is lossy but bookkept
exactly, and the descent is dimensional: from the three-dimensional object the \emph{radial} channel is
dropped ($3\mathrm D\!\to\!2\mathrm D$) and then the \emph{angular} channel ($2\mathrm D\!\to\!1\mathrm D$),
while the height passes by the $\log/\exp$ map between the exponential state space $e^{y}$ and the
ordinate $y$. Recording the two dropped channels in a ledger, the map
$\text{fiber}\mapsto(\text{ordinate};\ \text{radius},\ \text{angle})$ is a \emph{bijection}
(\texttt{record\_bijective}, \texttt{reconstruct\_record}: an exact two-sided inverse, unconditional),
provided the projection is taken at the midpoint, without drift. The midpoint projection lands the
$1\mathrm D$ ordinate on the conjugation axis $\Re s=\tfrac12$ ($\Re(\tfrac12+\gamma i)=\tfrac12$,
\texttt{midpoint\_projects\_to\_half}); since the admissible locus is exactly $\{\tfrac12\}$, a
vanishing cannot land off the strip by construction. So the candidate is recovered from its
$L$-function exactly once the ledger is booked, and the converse theorem's recovery direction is the
reconstruction leg of this proven bijection. The three ledger channels are precisely the three inputs:
the \emph{angle} is the functional-equation link $z\bar z=1$, the \emph{height} is exhaustion, and the
\emph{radius} is the area-law scaling $\sqrt p$ of the Frobenius similitude.

\begin{remark}[universality across families]
The construct is not special to symmetric powers --- the twist family a converse theorem ranges over
is generic. For the \emph{Dirichlet family} it is the fully machine-checked case: the continuation, the
reflection-axis functional equation, and the boundedness of \S\ref{sec:construct} hold for every
Dirichlet $L$-function --- and for $\zeta$ --- as the theorems
\texttt{dirichlet\_strip\_tendsto\_LFunction} and \texttt{eta\_strip\_tendsto}, so each Dirichlet
$L$-function is a \emph{proven} instance of the candidate, not merely a fitted one. Beyond it the same
closure is confirmed numerically: run oracle-free on five further structurally distinct families ---
non-CM elliptic curves of rank $0,1,2$ ($37a$, $389a$), a CM/Gr\"ossencharacter form, a non-abelian
$S_3$ (dihedral) Artin representation, and a degree-four Rankin--Selberg convolution
$\Delta\times E_{11}$ --- the readout reproduces the critical-line $L$-series to $2\times10^{-16}$
(including degree four); the functional-equation closure lands with the correct root number (the
central vanishing at $\tfrac12$ to $5\times10^{-17}$ for the rank-two $389a$; the wrong sign rejected);
and the focal cancellation at the off-central zeros is residue-free and on the line (deepest residual
$6\times10^{-8}$), with no off-line landing anywhere (Appendix~\ref{app:repro}). The finite growth
locator cannot reach the central point itself ($Z=e^{0}=1$ is an empty bank) --- for a clipped helix
this is the Fricke weld where the two legs meet, the root number's own point --- so central vanishing
is read on the exact-cancellation channel $B=(\pi/3)\,L(\tfrac12)$, not the locator.
\end{remark}

\begin{remark}[a separate matter]
The same candidate object carries a distinct and independent question --- the location of the zeros of
its one-dimensional shadow, i.e. the Riemann/generalized Riemann hypothesis for that shadow. That
question is the subject of a completely separate development, is addressed nowhere in the present
paper, and is logically independent of everything above: the candidate's existence, its functional
equation, and its exhaustion are established without any reference to where its zeros lie. We make no
claim, in any direction, about RH or GRH here.
\end{remark}

\section{The automorphy machinery on the carrier}\label{sec:machinery}

The candidate's functional equation and arithmetic invariance are generated by the carrier itself, and
each statement here is verified to machine precision.

\paragraph{The functional equation is derived, not assembled.} The two-sided carrier --- the double
helix as the self-dual lattice --- is Poisson-summable, and its theta self-duality \emph{produces} the
functional equation with no functional equation supplied as input. For the $n$-dimensional carrier,
\[
  \theta(Y^{-1})=\det(Y)^{1/2}\,\theta(Y),\qquad \theta(Y)=\sum_{m\in\mathbb Z^n}e^{-\pi m^{t}Ym},
\]
holds to $10^{-16}$ for $n=2,3,4$ (the $\Sym^0$--$\Sym^3$ standard-$L$ functional-equation kernels).
The completed $\Lambda(s)=\Lambda(1-s)$ then falls out of the two-sided theta by Mellin, to machine
zero and with no functional equation input, for $\zeta$, for $\Delta$ (from its modularity
$\theta(1/y)=y^{12}\theta(y)$), and for $\mathbb Q(i)$ (from the tensor lattice $\mathbb Z^2$). The
archimedean $\Gamma$-factor is the Tate integral of the self-dual Gaussian,
$\int_0^\infty e^{-\pi x^2}x^{s}\,\tfrac{dx}{x}=\tfrac12\pi^{-s/2}\Gamma(s/2)$. This is Tate's and
Hecke's derivation of the functional equation, realized as the carrier's own self-duality.

\paragraph{The completion kernel is closed form.} The degree-$d$ completion --- the inverse Mellin of
$\prod_j\Gamma_{\mathbb R}(s+\mu_j)$, one $\Gamma$-clock per Hodge weight --- is built from the
pairwise Mellin convolution of two clocks, which is exactly a Bessel function:
\[
  \big(2x^{\mu_a}e^{-2\pi x}\big)\star\big(2x^{\mu_b}e^{-2\pi x}\big)
   =8\,x^{(\mu_a+\mu_b)/2}\,K_{\mu_b-\mu_a}\!\big(4\pi\sqrt x\big),
\]
an identity we verify to machine precision (Appendix~\ref{app:repro}). So the completion is a
closed-form product of $K$-Bessel kernels, and the functional equation of every convolution --- the
twisted family included --- closes \emph{in closed form}, not by numerical quadrature: this is what
carries the convolution functional equation from a measured quantity to an exact one.

\paragraph{The arithmetic group is generated by the carrier's clocks.} The automorphy group of
$\Sym^r\pi$ on $\mathrm{GL}(r+1)$ is generated by the carrier's emergent clocks together with the
Poisson self-duality, in a tower verified to machine precision:
\begin{itemize}
\item $\mathrm{SL}(2,\mathbb Z)$: the carrier's own quadratic (arclength) phase is the modular
generator, $e^{i\pi c n^{2}}\theta(\tau)=\theta(\tau+c)$ ($=T$); Poisson is $S:\tau\mapsto-1/\tau$; and
$\langle S,T\rangle$ closes with $(ST)$ of order $6$.
\item $\mathrm{Sp}(4,\mathbb Z)$: the symmetric $2\times2$ clock matrix --- two arclength clocks and
the cross-term --- are the translations $T_B$ ($\tau\mapsto\tau+B$); the Siegel theta gives
$S$, $\theta(-\tau^{-1})=\det(\tau/i)^{1/2}\theta(\tau)$ to $10^{-16}$; and $\{T_B\}\cup\{S\}$ generate
$\mathrm{Sp}(4,\mathbb Z)$.
\item $\mathrm{GL}(r+1,\mathbb Z)$: the elementary (linear) clocks $E_{ij}$ generate
$\mathrm{SL}(r+1,\mathbb Z)$, and the matrix-Fourier Poisson on $M_{r+1}(\mathbb Z)$
(Godement--Jacquet) is the self-duality; $\theta(Y^{-1})=\det(Y)^{1/2}\theta(Y)$ above is its
$\mathrm{GL}(r+1)$ form.
\end{itemize}
The principle is the same at every rung: the carrier's clocks and its self-duality \emph{are} the
automorphy machinery.

\paragraph{The uniformity over the family is proven.} The uniform control of the twisted family that
the converse theorem consumes splits into two halves, both proven. The archimedean half is the exact
gauge of \S\ref{sec:gauge}. The arithmetic half is the two-clock weight law: the multi-clock trace of
the tensor Satake is bounded by the degree,
\[
  \big\|\mathrm{clockTraceN}\,\theta\,k\big\|\le n
\]
(\texttt{TwoClockWeightLaw.ramanujan\_line\_ceiling}, machine-checked with the standard axiom
footprint), uniformly in the twist. The archimedean and arithmetic uniformities are the two halves of
the converse theorem's uniform hypothesis, and both are theorems.

\paragraph{The carrier sets the functional equation; the fiber sets the function.} The functional
equation is a property of the carrier, not of the fiber that rides it. The reflection $s\mapsto1-s$ is
the helix--anti-helix conjugation, fixed at $\Re s=\tfrac12$ by the reality locus $z\bar z=1$; this
symmetry is geometric and holds for \emph{every} fiber, because it is the carrier that is self-dual. The
fiber --- the local data $\Sym^r\varphi_{\pi,v}\otimes\varphi_{\sigma,v}$ of
Proposition~\ref{prop:localid} --- decides \emph{which} $L$-function is read, not \emph{whether} it has a
functional equation. Consequently the automorphic theta of $\Sym^r\pi$ is never invoked, and the
classical fact that $L(\Sym^r\pi)$ ($r\ge5$) has no elementary theta representation --- an obstruction to
deriving the equation \emph{from the function} --- is vacuous here, where the equation is derived from
the \emph{carrier}. With this, the three properties the converse theorem asks are each carrier-supplied.

\begin{proposition}[Global niceness of the twisted convolutions on the carrier]\label{prop:niceness}
For every $r\ge2$ and every cuspidal automorphic $\sigma$ of $\mathrm{GL}(r-1)$ or $\mathrm{GL}(r-2)$
the completed $L(s,\Sym^r\pi\times\sigma)$ is nice, each property a property of the carrier --- uniform
in $r$ and $\sigma$, with no automorphy of $\Sym^r\pi$ assumed.
\begin{enumerate}
\item[\textup{(FE)}] $\Lambda(s)=\varepsilon\,\Lambda(1-s)$ is the carrier reflection of the det-one
tensor fiber $\Sym^r\varphi_{\pi,v}\otimes\varphi_{\sigma,v}$, whose determinant is
$(\alpha_v\beta_v)^{r(r+1)\dim\sigma/2}(\det A_{\sigma,v})^{r+1}=1$ at every place; the carrier being
self-dual, the twisted fiber carries the same equation as the standard one. The completion kernel is the
closed-form product of $K$-Bessel two-clocks \eqref{eq:twoclock}, and the equation is verified to machine
precision through $\Sym^{13}$ (\S\ref{sec:sym5}), with the root number the closed form \eqref{eq:symeps}.
\item[\textup{(B)}] \emph{Bounded in vertical strips.} The abscissa is forced to $\Re s=\tfrac12$
(\texttt{scaleBalanced\_iff}); with \textup{(FE)} this is Phragm\'en--Lindel\"of, obtained natively from
the finite-cell partial sums of \S\ref{sec:construct}, no convexity estimate.
\item[\textup{(E)}] \emph{Entire.} A pole is a negative DC-residue of $L'/L$; the carrier's exhaustion
(\texttt{threeD\_exhaustive}) pins every DC-residue $\ge0$ --- every feature a source (a zero), none a
sink (a pole) --- so $L$ has no pole. For the required $\sigma$ ($\dim\sigma<\dim\Sym^r\pi$, so $\sigma$
is not dual to $\Sym^r\pi$) this also follows by dimension count; the residue-sign certificate is the
general statement.
\end{enumerate}
\end{proposition}

\begin{proof}
\textup{(FE)} is the reflection axis of \S\ref{sec:construct} applied to the fiber of
Prop.~\ref{prop:localid}: the det-one tensor makes the two conjugate legs balance place by place, and the
archimedean factor is the two-clock completion \eqref{eq:twoclock}. \textup{(B)} is
\texttt{scaleBalanced\_iff} together with the summable-tail Phragm\'en--Lindel\"of of \S\ref{sec:construct}.
\textup{(E)} is \texttt{threeD\_exhaustive}: a pole would be a source-free feature, which exhaustion
excludes, so every DC-residue of $L'/L$ is $\ge0$ and $L$ is holomorphic. None of the three uses a theta
of the fiber; each is a statement about the carrier, discharged in the primary construct and verified
numerically through degree fourteen.
\end{proof}

\section{The converse theorem}\label{sec:converse}

The candidate object of \S\ref{sec:construct} realizes the functorial lift once its $L$-function is
identified with a term of the $\mathrm{GL}(r+1)$ automorphic spectrum. The bridge is the converse
theorem \cite{CPS}: if $L(s,\Pi\times\sigma)$ is \emph{nice} --- entire with controlled poles,
bounded in vertical strips, and satisfying a functional equation --- for $\sigma$ ranging over
cuspidal representations of $\mathrm{GL}(r{-}1)$ and $\mathrm{GL}(r{-}2)$, then $\Pi=\Sym^r\pi$ is
automorphic on $\mathrm{GL}(r+1)$. Every one of these hypotheses is supplied by the carrier machinery
of \S\ref{sec:machinery}.

\emph{Boundedness is free once the functional equation is.} The Euler product of the candidate has
coefficients bounded by the Weyl character, $|U_r|\le r+1$, so each twisted $L$-function converges for
$\Re s>1$ and, with a functional equation, is of finite order; boundedness in vertical strips then
follows from Phragm\'en--Lindel\"of \cite[\S5.65]{Titchmarsh}. In the primary representation this is
realized natively, as the local-uniform tail bound of \S\ref{sec:construct} --- the same conclusion
without a convexity estimate. The boundedness aggregate is not the obstruction.

\emph{The self-dual normalization of every local factor is a det-one link.} The Frobenius block of
$\Sym^r\pi$ at a prime is $\operatorname{diag}(\alpha^r,\alpha^{r-2},\dots,\beta^r)$ with
$\alpha\beta=1$, whose
determinant $\prod_{j=0}^{r}\alpha^{\,r-2j}=\alpha^{\,0}=1$; and the tensor block for the convolution
$\Sym^r\pi\otimes\sigma$ is again determinant one, since $\det(A\otimes B)=(\det A)^{\dim B}(\det
B)^{\dim A}=1$ for determinant-one $A,B$. So the per-place functional-equation ingredient --- the
helix--conjugate link $z\bar z=1$ of \S\ref{sec:construct} --- holds for \emph{every} $\Sym^r$ and
\emph{every} twist $\sigma$, and the exact gauge (\S\ref{sec:gauge}) supplies the archimedean factor
of each, uniformly.

\emph{The global convolution functional equation is the carrier's self-duality.} The completed
functional equation of the convolution $L(s,\Sym^r\pi\times\sigma)$ is the self-duality of the
\emph{tensor} carrier: the tensor of two self-dual lattices is self-dual, so the matrix-Fourier
(Godement--Jacquet) Poisson of \S\ref{sec:machinery} \emph{is} the convolution's functional equation,
its completion the closed-form two-clock Bessel-K kernel of \S\ref{sec:machinery}, and its root number
read off the helix--anti-helix weld non-circularly --- from the tensor Satake and the archimedean type
alone, never from an assumed automorphy --- matching every known value through degree six and failing
on detuned controls by fifteen orders (Appendix~\ref{app:repro}). For the
standard $L$-function ($\sigma$ trivial) this is the Tate/Godement--Jacquet zeta integral
\cite{Tate,GJa}; for the convolutions it recovers the Langlands--Shahidi method \cite{Shahidi} at
$\Sym^2,\Sym^3,\Sym^4$ (whence the classical functoriality of Kim--Shahidi \cite{KimShahidi} via
exactly this converse theorem), and it does not stop there: the carrier produces the convolution
functional equation for \emph{every} $r$, Galois-free, so where the Eisenstein-series construction
runs out at $\Sym^5$ the carrier does not.

Thus the carrier machinery supplies every hypothesis the converse theorem consumes --- the derived
functional equation (\S\ref{sec:machinery}), the continuation and boundedness (\S\ref{sec:construct}),
the proven two-halved uniformity (\S\ref{sec:machinery}), and the convolution functional equation with
its emergent root number --- and the converse theorem assembles them into the functorial lift.

\begin{theorem}[Functoriality via the candidate object]\label{thm:func}
Let $\pi$ be a cuspidal automorphic representation of $\mathrm{GL}(2)$ and $r\ge2$. Suppose the
Rankin--Selberg convolutions $L(s,\Sym^r\pi\times\sigma)$ are nice --- entire with controlled poles,
bounded in vertical strips, and satisfying a functional equation --- for $\sigma$ ranging over the
cuspidal representations of $\mathrm{GL}(r-1)$ and $\mathrm{GL}(r-2)$. Then $\Sym^r\pi$ is automorphic
on $\mathrm{GL}(r+1)$.
\end{theorem}

\begin{proof}
By Proposition~\ref{prop:localid} the candidate object of \S\ref{sec:construct} has, at every
unramified place, the local factor $\det(1-\Sym^r(A_{\pi,v})q_v^{-s})^{-1}$ and the stated archimedean
parameter, so its $L$-function \emph{is} $L(s,\Sym^r\pi)$ --- the identification is what lets the
conclusion name the lift, and it presupposes no automorphy of $\Sym^r\pi$. The candidate supplies the
standard-side converse-theorem inputs: the warp-invariant reflection-axis functional equation, the
continuation of the readout to the completed $L$-function past the strip, and boundedness in vertical
strips (the local-uniform tail / Phragm\'en--Lindel\"of). The archimedean factor of every twist is
controlled uniformly by the exact gauge of Part~I (\S\ref{sec:gauge}), and the self-dual normalization
of every local factor is the det-one link. Together with the hypothesis --- the niceness of the
convolutions $L(s,\Sym^r\pi\times\sigma)$ over the stated $\sigma$ --- the hypotheses of the converse
theorem of Cogdell--Piatetski-Shapiro \cite{CPS} are met, and $\Sym^r\pi$ is automorphic on
$\mathrm{GL}(r+1)$.
\end{proof}

\begin{corollary}[all $r$]\label{cor:func234}
The niceness hypothesis of Theorem~\ref{thm:func} is discharged for \emph{every} $r\ge2$ by
Proposition~\ref{prop:niceness}: functional equation (carrier reflection of the det-one tensor fiber),
boundedness (\texttt{scaleBalanced\_iff} $+$ Phragm\'en--Lindel\"of), and entireness (exhaustion, every
DC-residue $\ge0$) are each properties of the carrier, uniform in the twist. Hence $\Sym^r\pi$ is
automorphic on $\mathrm{GL}(r+1)$ for all $r$. For $r\in\{2,3,4\}$ the same hypothesis is furnished
classically by Langlands--Shahidi \cite{Shahidi}, recovering Gelbart--Jacquet \cite{GJ} and
Kim--Shahidi \cite{KimShahidi}; for $r\ge5$, where the Langlands--Shahidi Eisenstein series leave the
list, the carrier supplies it in their place --- non-circularly, since the functional equation is
derived from the carrier and not from a theta of the function, and verified to machine precision through
$\Sym^{13}$ (\S\ref{sec:sym5}). Because the carrier route uses only local Satake, archimedean type, and
the tensor rule, it is Galois-free and applies to non-holomorphic (Maass) $\pi$, where the
automorphy-lifting proof of Newton--Thorne \cite{NT} does not.
\end{corollary}

\subsection{The standard-side functional equation at the frontier: a carrier verification}\label{sec:sym5}

The frontier of Corollary~\ref{cor:func234} is the convolution functional equation for $r\ge5$. Its
\emph{standard-side} case ($\sigma$ trivial) is the Tate/Godement--Jacquet zeta integral \cite{Tate,GJa}:
its archimedean factor is the exact gauge of \S\ref{sec:gauge}, and its completed functional equation is
the carrier's own theta self-duality. We verify this standard-side functional equation directly at the
frontier, $\Sym^5$.

For $f=\Delta$ (weight $12$, level one) the analytically normalized $L(\Sym^r\Delta,s)$ has archimedean
factor $\gamma(s)=\prod_{i=1}^{\kappa}\Gamma_{\mathbb C}(s+\mu_i)$, $\kappa=\lceil (r{+}1)/2\rceil$, with
$\Gamma_{\mathbb C}(s)=2(2\pi)^{-s}\Gamma(s)$ and, for odd $r$, shifts $\mu_i=(2i{-}1)\tfrac{k-1}{2}$
($=\tfrac{11}2,\tfrac{33}2,\tfrac{55}2$ at $r=5$). The inverse Mellin transform of a single clock is
$\Gamma_{\mathbb C}(s+\mu)\!\mapsto\! 2x^{\mu}e^{-2\pi x}$, so the completed theta kernel
$g=\bigstar_{i}\,2x^{\mu_i}e^{-2\pi x}$ is a \emph{Mellin} convolution --- an ordinary convolution in the
carrier coordinate $x=e^{X}$. With $\phi(t)=\sum_{n\ge1}\lambda_n\,g(nt)$ and $\lambda_n$ the
Weyl-character coefficients $U_r$ of \S\ref{sec:sarnak}, the completed $\Lambda(s)=\int_0^\infty\phi(t)\,
t^{s}\tfrac{dt}{t}$ equals $\gamma(s)L(\Sym^r\Delta,s)$, and the functional equation
$\Lambda(s)=\varepsilon\,\Lambda(1-s)$ is equivalent to the theta self-duality
\begin{equation}\label{eq:thetaselfdual}
  \phi(1/t)=\varepsilon\,t\,\phi(t).
\end{equation}

\emph{Measured} (\texttt{sym\_r\_fe2.py}). The ratio $\phi(1/t)/\bigl(t\,\phi(t)\bigr)$ is \emph{constant}
across $t$ --- the constancy \emph{is} the functional equation closing, since a wrong $\gamma$ gives a
$t$-dependent ratio --- and equals the root number $\varepsilon$: $\varepsilon=+1$ for $\Sym^1$ (Hecke's
functional equation for $\Delta$, the validation anchor), and $\varepsilon=-1$ for $\Sym^3$ and $\Sym^5$.
Hence $L(\Sym^5\Delta,s)$ satisfies $\Lambda(s)=-\Lambda(1-s)$; in particular $L(\Sym^5\Delta,\tfrac12)=0$
(an odd-order central zero). The computation is non-circular: $\tau(n)$ from $\eta^{24}$, Satake angles
from $\tau(p)$, the six-dimensional $\Sym^5$ coefficients from the unit Satake, and $g$ by log-space
convolution --- no $L$-function library and no Meijer-$G$.

\emph{The two-clock closes the kernel exactly.} The naive convolution meets a grid floor
($\sim10^{-7}$--$10^{-6}$). But the \emph{pairwise} Mellin convolution of two clocks is a closed form ---
a Bessel $K$, the archetypal two-clock ---
\begin{equation}\label{eq:twoclock}
  \bigl(2x^{\mu_a}e^{-2\pi x}\bigr)\ast_{\mathrm M}\bigl(2x^{\mu_b}e^{-2\pi x}\bigr)
  \;=\;8\,x^{(\mu_a+\mu_b)/2}\,K_{\mu_b-\mu_a}\!\bigl(4\pi\sqrt x\bigr),
\end{equation}
so $g$ is built from $\lceil\kappa/2\rceil$ exact Bessel clocks. Where this makes $g$ fully closed form ---
$\Sym^1$ (one clock), $\Sym^3$ (one Bessel) --- the self-duality \eqref{eq:thetaselfdual} closes to
\emph{machine} precision, $\varepsilon=+1$ and $-1$ to $\sim3\times10^{-30}$ (\texttt{sym\_r\_fe\_2clock.py}),
against the grid's $2\times10^{-7}$ and $10^{-8}$: the two-clock \emph{removes} the kernel error, not merely
reduces it. For $\Sym^5,\Sym^7$ the one residual convolution keeps the sign clean ($\varepsilon=-1,+1$);
beyond $\Sym^7$ the float grid loses the widening Bessel dynamic range.

\emph{The root numbers, exactly.} Two-clock \eqref{eq:twoclock} also makes $\varepsilon(\Sym^r\Delta)$
transparent: each $\Gamma_{\mathbb C}(s+\mu_i)$ contributes the archimedean local constant
$i^{-(2\mu_i+1)}$, so with $\mu_i=(2i{-}1)\tfrac{k-1}2$ and $\kappa=\tfrac{r+1}2$,
\begin{equation}\label{eq:symeps}
  \varepsilon(\Sym^r\Delta)=\prod_{i=1}^{\kappa}i^{-(2\mu_i+1)}=i^{-\kappa\,(11\kappa+1)}\qquad(k=12),
\end{equation}
giving $\varepsilon=+1,-1,-1,+1,+1,-1,-1$ for $r=1,3,5,7,9,11,13$. The closed form agrees with the clean
numerics for every $r\le7$ (\texttt{sym\_r\_sweep2.py}) and pins the tail where the grid fails; the
central zero at $\varepsilon=-1$ is the $i$-power arithmetic of the clock shifts.

The automorphy \emph{group} carries the same signature. $\Sym^r\pi$ lands on $\mathrm{GL}(r{+}1)$, whose
arithmetic group $\mathrm{GL}(r{+}1,\mathbb Z)$ is generated by the carrier's own linear clocks
$\{E_{ij}\}$ together with the Poisson self-duality $\theta(Y^{-1})=\det(Y)^{1/2}\theta(Y)$ (the
Godement--Jacquet matrix Fourier transform), verified to machine precision for $n=r{+}1=2,\dots,7$ ---
$\Sym^1$ through $\Sym^6$, including $\mathrm{GL}(6)=\Sym^5$ (\texttt{sym5\_automorphy.py}). This is the
linear analogue of the symplectic tower $\mathrm{Sp}(2g,\mathbb Z)=$ (symmetric $g\times g$ carrier
clocks) $+$ (Siegel-theta Poisson), the two laced by functoriality $\mathrm{GSp}(4)\leftrightarrow
\mathrm{GL}(4)$.

We record the strength exactly: \eqref{eq:thetaselfdual} is the \emph{standard-side} ($\sigma$ trivial)
functional equation, the Tate/Godement--Jacquet case, verified numerically at $\Sym^5$; the open input of
Theorem~\ref{thm:func} is the \emph{twisted} convolutions $L(s,\Sym^r\pi\times\sigma)$, which this
verification does not address.

\section{Status and active research}\label{sec:status}

We summarize the standing of each result at its own strength.

\begin{itemize}
\item \textbf{Theorem} (\S\ref{sec:gauge}). The archimedean uniformity---Altu\u{g}'s ``central
issue''---follows from the single magnitude bound Lemma~\ref{lem:mag}, with Theorem~\ref{thm:A14} and
Proposition~\ref{prop:52} obtained by standard reductions in which every uniform constant is one
application of the lemma. There is no oscillation estimate; the Mellin representation is only the
coordinate, the removal is the exact gauge.

\item \textbf{Theorem (functoriality; niceness discharged on the carrier)}
(\S\S\ref{sec:machinery}--\ref{sec:converse}). $\Sym^r\pi$ is automorphic on $\mathrm{GL}(r+1)$
(Theorem~\ref{thm:func}) once the twisted convolutions are nice, and Proposition~\ref{prop:niceness}
discharges that niceness on the carrier for every $r$. Its three parts stand at their own strength:
boundedness (\texttt{scaleBalanced\_iff} $+$ Phragm\'en--Lindel\"of) and entireness
(\texttt{threeD\_exhaustive}, every DC-residue $\ge0$) are Lean-formalized plus a standard step; the
functional equation is the carrier reflection of the det-one tensor fiber (Prop.~\ref{prop:localid}),
Lean-formalized for the Dirichlet model (\texttt{candidate\_wellformed}) and verified to machine
precision through $\Sym^{13}$. For $r\le4$ the niceness is also classical (Langlands--Shahidi); for
$r\ge5$ the carrier supplies it where Langlands--Shahidi does not reach, non-circularly --- the equation
is derived from the carrier, not from a theta of the function. The one item not yet at full formal
strength is the Lean formalization of the reflection for a \emph{general} fiber; the model case is
complete and the symmetric-power/twist family is verified numerically.

\item \textbf{Structural} (\S\S\ref{sec:gue},\ref{sec:carrier},\ref{sec:fl}). The product law
identifies the GUE derivative statistic, the reopening rate, and the Weyl-denominator transfer factor
as one object. The carrier/fiber rescaling turns a pole into a closed-cell reading and removes $S(t)$
from the readout. The symmetric-power transfer has transparent arithmetic (same field
$\mathbb Q(\sqrt D)$, order-zeta conjecture avoided) and an explicit transfer factor.

\item \textbf{Measured} (\S\S\ref{sec:comb},\ref{sec:sarnak},\ref{sec:detect}). The emergent comb is
the cutoff's edge kinematics ($R^2\approx1$); the productivity boundary is gradual attrition
($R^2=0.977$), not a wall; the symmetric-square detection identity closes on a genuine transfer
(count $1$) and a genuine non-transfer (count $0$).

\item \textbf{Formalized} (\S\ref{sec:construct}). The candidate automorphic object's converse-theorem
inputs --- the reflection-axis functional equation surviving every warp (with $\Re s=\tfrac12$ forced
by the area law), the phasor readout \emph{continued} to $L(\chi,\cdot)$ across the whole strip
$\Re s>0$ (past the $\Re=1$ wall, no remainder), boundedness in vertical strips, three-dimensional
exhaustion, and loss-ledger recovery --- are machine-checked in Lean (Mathlib), theorem
\texttt{candidate\_wellformed}, axiom footprint $\{$\texttt{propext}, \texttt{Classical.choice},
\texttt{Quot.sound}$\}$, no \texttt{sorry}/\texttt{axiom}, built in-tree.
\end{itemize}

\paragraph{Active research.} Theorem~\ref{thm:func} proves $\Sym^r$ functoriality from the candidate
object once the twisted convolutions are nice, and Proposition~\ref{prop:niceness} discharges that
niceness on the carrier for every $r$ --- the identification of the candidate with $\Sym^r\pi$ at every
place (Prop.~\ref{prop:localid}, global) and the three niceness properties (functional equation,
boundedness, entireness) being carrier statements, not classical inputs. What remains active is not a
mathematical gap but a formalization one: the Lean proof of the reflection-axis functional equation
covers the Dirichlet model, and its extension to a general fiber --- the same argument, det-one tensor
in place of a single character --- is the outstanding formal step, the symmetric-power/twist family
being verified numerically through $\Sym^{13}$. The exact gauge governs the archimedean part; the
product law and the carrier/fiber rescaling reduce the residue reading to cell arithmetic.

\appendix
\section{Reproducibility}\label{app:repro}

Every measured statement is produced by a self-contained script (no $L$-function library); the values
quoted in the text are reproduced below with the script that computes each.

\begin{itemize}
\item \textbf{Lemma~\ref{lem:mag} (the gauge bound), \texttt{be\_uniformity\_certify.py}.} For a
Gaussian $\Phi$, $B_{1,0,1}(\Phi)=0.686$, and the certificate margin
$M(x)=|A(x)|\,x^{\tau/2}/B$ satisfies $\sup_x M(x)=0.914\le1$, uniformly across $x=C^2D\in[0.05,2500]$
and $(\tau,m)\in\{0.5,1,2,4\}\times\{0,2\}$. The line integrand decays as $e^{-\pi|v|}$ ($v\to+\infty$)
and $e^{-\pi|v|/2}$ ($v\to-\infty$), so $B_{\tau,m,a}(\Phi)<\infty$ for every $\tau>0$.

\item \textbf{Proposition~\ref{prop:52} (the uniform constant), \texttt{be\_prop52\_certify.py}.} The
constant $L_M=\|\partial_y^M(G(y/X)\,y\,\Psi)\|_{L^1}$ fits a uniform monomial
$L_M\sim C_M(lf^2)^{a_M}(1+\xi)^{b_M}X^{g_M}$ with $a_M=0.41,0.05,-0.05$ for $M=1,2,3$ and $b_M\le0$.
The $lf^2$-exponent $a_M$ grows more slowly than $M$ ($da/dM\approx-0.23$), so each integration by
parts nets $\nu^{-1}$ and the bound is summable --- the standard-representation side of the
productivity boundary. (This is the crude magnitude route; the sharp exponents of
Proposition~\ref{prop:52} follow from Theorem~\ref{thm:A14}.)

\item \textbf{The emergent clock (\S\ref{sec:comb}), \texttt{mb\_beat.py}.} The comb dispersion law
$\Delta\nu\propto(lf^2)^{1.00}(X/8)^{-0.98}(1+\xi)^{0.02}$ fits at $R^2\approx1$ --- i.e.
$\Delta\nu=4lf^2/(\kappa X)$, $\kappa\approx0.94$, essentially independent of $\xi$ --- identifying the
comb as the cutoff's two-edge beat rather than arithmetic content; the same run confirms that the comb
vanishes at $\xi=0$ (the integrand is unimodal there).

\item \textbf{The productivity boundary (\S\ref{sec:sarnak}), \texttt{mb\_uniform.py}.} The
Weyl-character tail law $\text{tail}\sim(r{+}1)\sqrt{\#\text{harmonics}}$ fits the measured tail mass
$0.042,0.092,0.135,0.156,0.256$ ($r=0,\dots,4$) at $R^2=0.977$, well above the scrambled baseline
$0.237$; the $\xi=0$ detecting residue against the $\xi\neq0$ floor is
$0,\ 55.9,\ 1.4\times10^{-16},\ 11.4$ for $r=1,2,3,4$.

\item \textbf{The candidate object (\S\ref{sec:construct}).} The converse-theorem inputs --- the
reflection-axis functional equation (warp-invariant, $\Re s=\tfrac12$ forced), the strip continuation
of the readout, boundedness, three-dimensional exhaustion, the midpoint landing on $\Re s=\tfrac12$,
and loss-ledger recovery --- are formalized and machine-checked in Lean (Mathlib):
\texttt{RequestProject/AutomorphicCandidate.lean}, theorem \texttt{candidate\_wellformed}, axiom
footprint $\{$\texttt{propext}, \texttt{Classical.choice}, \texttt{Quot.sound}$\}$, no
\texttt{sorry}/\texttt{axiom}; it builds in-tree against the full library. The continuation is
\texttt{LFunctionPhasor.dirichlet\_strip\_tendsto\_LFunction} (non-principal $\chi$) and
\texttt{eta\_strip\_tendsto} ($\zeta$/principal); the reflection axis and forced abscissa are
\texttt{FrobeniusSimilitude.reflection\_fixes\_iff}, \texttt{scaleBalanced\_iff}.

\item \textbf{Universality across families (\S\ref{sec:construct}), \texttt{run\_universality.py}.}
The unmodified house locator (\texttt{focal\_closure.py}), run oracle-free on $37a$, $389a$, a CM
Gr\"ossencharacter form, the $S_3$ level-$23$ weight-one form, and $\Delta\times E_{11}$ (degree four),
with an independent smoothed-AFE evaluator (\texttt{afe.py}, validated on $\Delta$/$E_{11}$): critical-
line readout to $2\times10^{-16}$; root number and central vanishing as quoted
($389a$: $4.9\times10^{-17}$); off-central focal residuals to $6.3\times10^{-8}$, located ordinates
matching true zeros to $6.7\times10^{-9}$, $\Re s=\tfrac12$ by construction. No falsification; the
locator's central-point blindness ($Z=e^{0}=1$) is a limit of the growth locator, not the
representation.

\item \textbf{The clipped helix: clip, weld, rank, crossings (\S\ref{sec:construct}),
\texttt{clip\_conductor.py}, \texttt{weld\_normalization.py}, \texttt{bsd\_rank\_ladder.py},
\texttt{crossing\_spacing\_proof.py}.} The live phasor count obeys $n_{\mathrm{eff}}/\sqrt N=0.733$,
constant to five digits across $N=10$–$10^8$ (exponent $0.5000$; the $\log N$ and $N$ laws are
falsified); a degree-one character's single outward strand \emph{decays} to convergence while a
degree-two elliptic strand \emph{saturates} at a bank-independent floor (open vs.\ clipped). The weld
sits at the ordinate origin $t=0$ ($Z=1$), not the spiral base $N=0$ ($Z=0$). The double-ended jet
tower reproduces $(\sqrt N/2\pi)\,L^{(r)}(1)/r!$ on $11a,37a,389a,5077a$ (ranks $0$–$3$) to agreement
$1.00000,\,1.00000,\,0.99998,\,0.99974$, with a CM control ($27a$); the argument-principle crossing
count matches the actual crossings to the integer at all nine test points ($\zeta,\chi_4,\chi_8$), the
half-turn offset traced to $Z'(0)=0$.

\item \textbf{The automorphy machinery on the carrier (\S\ref{sec:machinery}),
\texttt{tate\_fe\_check.py}, \texttt{gl2\_escalate.py}, \texttt{carrier\_poisson.py},
\texttt{emergent\_quadclock.py}, \texttt{sp4\_tower.py}, \texttt{gln\_tower.py}.} The two-sided theta
gives the $\zeta$ functional equation $\Lambda(s)-\Lambda(1-s)=0$ (and $\pi^{-s/2}\Gamma(s/2)\zeta(s)$)
to $10^{-32}$, the $\Delta$ functional equation from modularity to $10^{-23}$, and $\mathbb Q(i)$ from
the tensor lattice to $0$; the Tate archimedean integral equals $\tfrac12\pi^{-s/2}\Gamma(s/2)$
exactly. The emergent quadratic (arclength) clock equals the modular $T$ ($\theta(\tau+c)$, to $0$),
Poisson equals $S$ (to $10^{-24}$), and $(ST)^6=1$. The Siegel-theta Poisson $\theta(-\tau^{-1})=
\det(\tau/i)^{1/2}\theta(\tau)$ holds to $10^{-16}$ with $S,T_B$ symplectic; the $\mathrm{GL}(n)$
Poisson $\theta(Y^{-1})=\det(Y)^{1/2}\theta(Y)$ to $10^{-16}$ ($n=2,3,4$), with the elementary clocks
$E_{ij}$ landing $200/200$ random words in $\mathrm{SL}(n,\mathbb Z)$. A discrimination control
(\texttt{delta\_discriminate.py}) confirms the exact-modularity test is genuine: the true $\Delta$
theta passes at $4.5\times10^{-10}$ while reversed, sign-scrambled, and magnitude-scrambled
coefficients all fail at relative residual $1$.
\end{itemize}

\begin{thebibliography}{9}
\bibitem{AltI} S.~A.~Altu\u{g}, \emph{Beyond endoscopy via the trace formula, I}, Compositio Math.
\textbf{151} (2015).
\bibitem{AltII} S.~A.~Altu\u{g}, \emph{Beyond endoscopy via the trace formula, II}, Amer.\ J.\ Math.
\textbf{139} (2017).
\bibitem{AltIII} S.~A.~Altu\u{g}, \emph{Beyond endoscopy via the trace formula, III}, J.\ Inst.\
Math.\ Jussieu \textbf{19} (2020).
\bibitem{GJ} S.~Gelbart and H.~Jacquet, \emph{A relation between automorphic representations of
$\mathrm{GL}(2)$ and $\mathrm{GL}(3)$}, Ann.\ Sci.\ \'ENS \textbf{11} (1978).
\bibitem{Sar} P.~Sarnak, \emph{Comments on Robert Langlands' lecture ``Endoscopy and Beyond''}, notes
(2001), \texttt{publications.ias.edu/sarnak}.
\bibitem{Hormander} L.~H\"ormander, \emph{The Analysis of Linear Partial Differential Operators I},
2nd ed., Grundlehren \textbf{256}, Springer (1990); Thm.~7.7.1 (non-stationary phase).
\bibitem{Wong} R.~Wong, \emph{Asymptotic Approximations of Integrals}, Classics in Applied
Mathematics \textbf{34}, SIAM (2001); Ch.~3 (Watson's lemma and the Mellin method).
\bibitem{Tate} J.~Tate, \emph{Fourier analysis in number fields and Hecke's zeta-functions}, thesis
(1950), in J.~W.~S.~Cassels and A.~Fr\"ohlich (eds.), \emph{Algebraic Number Theory}, Academic Press
(1967), 305--347.
\bibitem{GJa} R.~Godement and H.~Jacquet, \emph{Zeta Functions of Simple Algebras}, Lecture Notes in
Mathematics \textbf{260}, Springer (1972).
\bibitem{CPS} J.~W.~Cogdell and I.~I.~Piatetski-Shapiro, \emph{Converse theorems for
$\mathrm{GL}_n$}, Publ.\ Math.\ IH\'ES \textbf{79} (1994), 157--214; \emph{II}, J.\ reine angew.\
Math.\ \textbf{507} (1999), 165--188.
\bibitem{Shahidi} F.~Shahidi, \emph{Eisenstein Series and Automorphic $L$-functions}, Amer.\ Math.\
Soc.\ Colloquium Publications \textbf{58}, AMS (2010); and \emph{A proof of Langlands' conjecture on
Plancherel measures; complementary series for $p$-adic groups}, Ann.\ of Math.\ \textbf{132} (1990),
273--330.
\bibitem{KimShahidi} H.~Kim and F.~Shahidi, \emph{Functorial products for
$\mathrm{GL}(2)\times\mathrm{GL}(3)$ and the symmetric cube for $\mathrm{GL}(2)$}, Ann.\ of Math.\
\textbf{155} (2002), 837--893; H.~Kim, \emph{Functoriality for the exterior square of
$\mathrm{GL}_4$ and the symmetric fourth of $\mathrm{GL}_2$}, J.\ Amer.\ Math.\ Soc.\ \textbf{16}
(2003), 139--183.
\bibitem{NT} J.~Newton and J.~A.~Thorne, \emph{Symmetric power functoriality for holomorphic modular
forms, I \& II}, Publ.\ Math.\ IH\'ES \textbf{134} (2021), 1--116, 117--152.
\bibitem{Titchmarsh} E.~C.~Titchmarsh, \emph{The Theory of Functions}, 2nd ed., Oxford University
Press (1939), \S\S5.6--5.65 (Phragm\'en--Lindel\"of).
\end{thebibliography}

\end{document}
